\documentclass[11pt,a4paper]{article}
\usepackage[margin=2.5cm]{geometry}
\usepackage[T1]{fontenc}
\usepackage{lmodern}
\usepackage{microtype}
\usepackage{amsmath,amssymb,amsthm}
\usepackage{booktabs}
\usepackage{array}
\usepackage{enumitem}
\usepackage{xcolor}
\usepackage{fancyvrb}
\usepackage{tikz-cd}
\usepackage[hidelinks]{hyperref}

\definecolor{codebg}{gray}{0.95}
\DefineVerbatimEnvironment{code}{Verbatim}{fontsize=\small,frame=single,framerule=0.2pt,framesep=2mm}

\newcommand{\level}[1]{\texttt{#1}}
\newcommand{\field}[1]{\texttt{#1}}

\theoremstyle{plain}
\newtheorem{theorem}{Theorem}[section]
\newtheorem{lemma}[theorem]{Lemma}
\newtheorem{corollary}[theorem]{Corollary}
\newtheorem{proposition}[theorem]{Proposition}
\theoremstyle{definition}
\newtheorem{definition}[theorem]{Definition}
\newtheorem{example}[theorem]{Example}
\newtheorem{remark}[theorem]{Remark}

\title{Trust Systems for Autonomous Agent Composition:\\
\large Theory, Formalization, and Deployment}
\author{Robin Langer \and Claude}
\date{July 2026}

\begin{document}
\maketitle

\begin{abstract}
We describe a trust system for autonomous research agents that grades
mathematical claims, validates citations, and composes across agent
boundaries via email. Three agents---Clio (algebraic combinatorics),
Rick (combinatorial Hopf algebras), and MacBeth (category theory)---each
run independent trust systems inside Docker containers with different
trust chains, different boundaries, and different evidence duties. A
monotone interface map $\varphi$ translates grades between systems;
a new grade \level{peer-claimed} mediates cross-agent citation via
email with proof attachments. The core theory (the certificate lemma,
trust transport, consolidation preservation) is Lean-verified with no
\texttt{sorry}. The registry structure is a directed container whose
positions separate into a proof tree (\field{role: "premise"}) and an
exploration tree (\field{role: "attempt"}). We describe the complete
system: definitions, theorems, Lean formalization, the generic
validator, and the operational deployment across four agents.
\end{abstract}

\tableofcontents

% ===================================================================
\section{Introduction}
% ===================================================================

An autonomous research agent proves theorems, reads papers, and
consolidates memory across cycles. Each of these operations can corrupt
knowledge: a sub-agent hallucinates paper contents, compression drops
citations, a proof absorbs hearsay as fact. The trust system exists to
make such corruption mechanically visible.

This paper describes the complete system as deployed across four agents.
It is simultaneously a mathematical paper (definitions, theorems, proofs),
a systems paper (the validator, the descriptors, the email protocol), and
a field report (what happened when the agents ran). The three contributions
are:

\begin{enumerate}
\item \textbf{The theory.} Trust systems $({\mathcal D}, T, s)$ with the
certificate lemma (Lemma~\ref{lem:certificate}), trust transport
(Lemma~\ref{lem:transport}), and consolidation preservation
(Theorem~\ref{thm:consolidation}). All Lean-verified.

\item \textbf{The composition mechanism.} Cross-agent citation via a
monotone interface map $\varphi$, mediated by email with proof
attachments, graded at \level{peer-claimed}. The first live composition
event: Rick peer-reviews Clio's claims; Clio upgrades nodes to
\level{peer-reviewed} with the review artifact.

\item \textbf{The tree pair.} Every child node carries a \field{role}
field (\field{"premise"} or \field{"attempt"}) that separates the
downward proof tree from the outward exploration tree. The boundary rule
checks only premises. This makes the directed container structure
explicit: the proof tree is the successful-path coKleisli morphism; the
exploration tree is the dead-ends morphism.
\end{enumerate}

% ===================================================================
\section{The data of a trust system}
\label{sec:trust-system}
% ===================================================================

\begin{definition}[Trust system]
\label{def:trust-system}
A \emph{trust system} is a triple $({\mathcal D}, T, s)$:
\begin{enumerate}
\item \textbf{Dependency structure ${\mathcal D}$}: a multicategory.
Objects are \emph{claims}. A multimorphism $f : y_1, \ldots, y_n \to x$
is a \emph{checking step}: a derivation of $x$ from premises
$y_1, \ldots, y_n$. The $0$-ary case $f : () \to x$ is a \emph{ground
check} (reading a source, a Lean compilation, a reproduced computation).

\item \textbf{Check-poset $T$}: a bounded chain (linear order with
$\bot$), ordered by strength of check. Meets and joins are min and max.

\item \textbf{Grading} $s : \mathrm{Mor}({\mathcal D}) \to T$, assigning
to each checking step the strength of the check it received.
\end{enumerate}
\end{definition}

$T$ is the \emph{deployment parameter}. Each agent instantiates it
differently:

\begin{center}
\begin{tabular}{lll}
\toprule
Agent & Chain $T$ & Boundary \\
\midrule
Clio & \level{speculative} $<$ \level{computed} $<$ \level{peer-claimed} $<$ \level{proved} & \level{proved} \\
& \quad $<$ \level{peer-reviewed} $<$ \level{published} $<$ \level{lean-verified} & \\
Rick & \level{hunch} $<$ \level{sketched} $<$ \level{computed} $<$ \level{peer-claimed} & \level{checked-sober} \\
& \quad $<$ \level{checked-sober} $<$ \level{proved} $<$ \level{lean-verified} & \\
MacBeth & (same as Clio) & \level{proved} \\
\bottomrule
\end{tabular}
\end{center}

\begin{remark}[Non-isomorphic chains]
Rick's chain has seven levels; Clio's and MacBeth's have seven (with
\level{peer-claimed} added between \level{computed} and \level{proved}).
Rick's chain includes \level{hunch}, \level{sketched}, and
\level{checked-sober}, which have no direct analogues in Clio's system.
The interface map $\varphi$ (Definition~\ref{def:morphism}) translates
between them.
\end{remark}

\begin{definition}[Grade of a derivation, canonical trust]
\label{def:canonical-trust}
The grade of a derivation $\delta = (f; \delta_1, \ldots, \delta_n)$ is
the weakest link:
\[
g(\delta) = s(f) \wedge \bigwedge_i g(\delta_i).
\]
The \emph{canonical trust} of a claim $x$ is the best derivation:
\[
\tau(x) = \bigvee \{ g(\delta) : \delta \text{ a derivation of } x \},
\quad \text{with } \bigvee \emptyset = \bot.
\]
\end{definition}

This is a provenance-semiring instance (Green--Karvounarakis--Tannen,
PODS 2007): the min-max specialization. The semiring is not the
contribution; the deployment is.

\textbf{Standing hypotheses.} ${\mathcal D}$ is finite and well-founded:
every claim's justification bottoms out in ground checks. Finiteness
makes all sups attained; well-foundedness is load-bearing
(Remark~\ref{rem:acyclicity}).

% ===================================================================
\section{The certificate lemma}
\label{sec:certificate}
% ===================================================================

\begin{definition}[Sound assignment]
\label{def:sound}
An assignment $t : \mathrm{Ob}({\mathcal D}) \to T$ is \emph{sound} if
for every claim $x$ with $t(x) > \bot$ there exists a checking step
$f : y_1, \ldots, y_n \to x$ with
\begin{equation}
\label{eq:dagger}
t(x) \leq s(f) \wedge \bigwedge_i t(y_i). \tag{$\dagger$}
\end{equation}
\end{definition}

This is the \textbf{boundary rule} as deployed: a registry node may
claim level $\ell$ only if the step-check reached $\ell$ and every
premise child is at least $\ell$. Condition~\eqref{eq:dagger} is
node-local---checkable without seeing the whole tree. The validator is a
certificate checker: it verifies~\eqref{eq:dagger} at each node; it
never computes $\tau$.

\begin{lemma}[Certificate lemma]
\label{lem:certificate}
Under the standing hypotheses:
\begin{enumerate}[label=(\roman*)]
\item $\tau$ is sound.
\item Every sound $t$ satisfies $t \leq \tau$ pointwise.
\end{enumerate}
Hence $\tau$ is the greatest sound assignment.
\end{lemma}

\begin{proof}
(i) Suppose $\tau(x) > \bot$. Since $T$ is a finite chain, the sup is
attained: $\tau(x) = g(\delta)$ for some $\delta = (f; \delta_1, \ldots,
\delta_n)$. Then $\tau(x) = s(f) \wedge \bigwedge_i g(\delta_i) \leq
s(f) \wedge \bigwedge_i \tau(y_i)$, so~\eqref{eq:dagger} holds at $x$.

(ii) By well-founded induction. Let $t$ be sound and pick $f$ as
in~\eqref{eq:dagger}. By induction $t(y_i) \leq \tau(y_i)$; choose
$\delta_i$ with $g(\delta_i) = \tau(y_i)$. Then $\delta = (f; \delta_1,
\ldots, \delta_n)$ has $g(\delta) \geq t(x)$, so $\tau(x) \geq t(x)$.

\emph{Lean:} \texttt{sound\_le\_tau}, \texttt{tau\_sound} in
\texttt{Trust.lean}, no sorry.
\end{proof}

\begin{remark}[Well-foundedness is load-bearing]
\label{rem:acyclicity}
Drop well-foundedness and (ii) fails: two claims citing each other, both
graded $\top$, have $t \equiv \top$ sound but $\tau \equiv \bot$.
Circular citation is trust inflation that passes every node-local check.
The validator checks acyclicity explicitly via DFS across agent
namespaces.

\emph{Lean:} \texttt{cyc\_sound\_top}, \texttt{cyc\_no\_deriv},
\texttt{cyc\_not\_wf}.
\end{remark}

% ===================================================================
\section{Morphisms and cross-agent citation}
\label{sec:morphisms}
% ===================================================================

\begin{definition}[Morphism of trust systems]
\label{def:morphism}
A \emph{morphism} $(F, \varphi) : ({\mathcal D}, T, s) \to ({\mathcal D}',
T', s')$ is a multifunctor $F : {\mathcal D} \to {\mathcal D}'$ together
with a monotone map $\varphi : T \to T'$ that is \textbf{strict at bottom}
($\varphi(\bot) = \bot'$) and laxly compatible with grading:
\[
\varphi(s(f)) \leq s'(Ff) \quad \text{for every checking step } f.
\]
\end{definition}

Strictness at $\bot$---``a morphism maps \emph{unchecked} to
\emph{unchecked}''---was found missing by the Lean formalization: without
it, transport fails on a claim with no derivation.

\begin{lemma}[Transport of trust]
\label{lem:transport}
For any morphism $(F, \varphi)$ and claim $x$:
\[
\varphi(\tau(x)) \leq \tau'(Fx).
\]
\end{lemma}

\begin{proof}
Induction on derivations. For $\delta = (f; \delta_1, \ldots, \delta_n)$,
$F\delta$ is a derivation of $Fx$ with
$g'(F\delta) \geq \varphi(g(\delta))$ by monotonicity and lax
compatibility. Taking sups and using attainment (finite chain) or
strictness at $\bot$ (empty derivation set).

\emph{Lean:} \texttt{tau\_transport} in \texttt{Trust.lean}, no sorry.
Both case splits found by the formalization.
\end{proof}

\begin{corollary}[Cross-agent citation]
\label{cor:cross-agent}
Clio citing a Rick registry node is composition of morphisms: Rick's
system $\to$ opaque restriction at his registry roots $\to$ Clio's system
along $\varphi : T_{\text{Rick}} \to T_{\text{Clio}}$.
Lemma~\ref{lem:transport} composed twice gives the safe import grade.
\end{corollary}

\subsection{The interface maps}

The $\varphi$ maps are declared in JSON descriptors. The critical rows:

\begin{center}
\begin{tabular}{lll}
\toprule
Direction & Key mapping & Effect \\
\midrule
Rick $\to$ Clio & \level{checked-sober} $\mapsto$ \level{peer-claimed} & Below Clio's boundary \\
Rick $\to$ Clio & \level{proved} $\mapsto$ \level{proved} & At Clio's boundary \\
Clio $\to$ Rick & \level{peer-claimed} $\mapsto$ \level{peer-claimed} & Below Rick's boundary \\
Clio $\leftrightarrow$ MacBeth & identity on all grades & Same chain \\
\bottomrule
\end{tabular}
\end{center}

Rick's \level{checked-sober} maps to \level{peer-claimed} (not
\level{computed}) in Clio's system. This is forced by monotonicity:
\level{peer-claimed} sits between \level{computed} and \level{proved}
in both chains, and the mapping must preserve order.

\subsection{Source-extraction demotion}

Each interface also declares a \texttt{source\_phi} that demotes
extraction levels by one step: deep-read by proxy becomes abstract-level
check locally. A sub-agent's summary of a paper that Rick read is two
degrees of separation from the mathematics.

\subsection{The validator checks}

The generic validator (\texttt{trustcheck.py}) enforces:
\begin{enumerate}
\item \textbf{$\varphi$ side conditions}: total, monotone, strict at $\bot$.
\item \textbf{Trust cap}: a shared stub's local trust $\leq \varphi(\text{canonical trust})$.
\item \textbf{Artifact transport}: if $\varphi(\text{canonical trust}) \geq$ boundary and the canonical node has a proof file, the file must exist under the peer's artifact directory.
\item \textbf{Acyclicity}: DFS across agent namespaces detects circular citation.
\end{enumerate}

% ===================================================================
\section{Consolidation preservation}
\label{sec:consolidation}
% ===================================================================

\begin{definition}[Consolidation]
\label{def:consolidation}
A \emph{consolidation} $C : ({\mathcal D}, T, s) \to ({\mathcal D}', T, s')$
over the same check-chain consists of:
\begin{itemize}
\item \textbf{Shapes forward}: a map on claims $x \mapsto Cx$ and on
checking steps $f \mapsto Cf$ with premises corresponding and grades not
dropping: $s(f) \leq s'(Cf)$.
\item \textbf{Positions backward}: for every claim $x$ and every checking
step $f' : y'_1, \ldots, y'_m \to Cx$ of the target, a chosen step
$\rho(f')$ of the source with premises corresponding and the target grade
cached: $s'(f') \leq s(\rho(f'))$.
\end{itemize}
\end{definition}

$\rho$ is the position map; a citation is precisely its cached value at
one step. The definition never mentions $\tau$: it is a bisimulation-style,
node-local condition.

\begin{theorem}[Consolidation preservation]
\label{thm:consolidation}
If $C$ is a consolidation and both premise relations are well-founded,
then $\tau'(Cx) = \tau(x)$ for every source claim $x$.
\end{theorem}

\begin{proof}
Forward: Lemma~\ref{lem:transport} at $\varphi = \mathrm{id}$. Backward:
well-founded induction on the source, using attainment, the premise
correspondence, and the cached-grade condition.

\emph{Lean:} \texttt{tau\_consolidate}, no sorry.
\end{proof}

\begin{theorem}[Consolidations compose]
\label{thm:consolidation-comp}
Consolidations have identities and compose: shapes forward through both
factors, positions backward (contravariantly), caches by transitivity.

\emph{Lean:} \texttt{Consolidation.id}, \texttt{Consolidation.comp},
\texttt{tau\_consolidate\_comp}, no sorry.
\end{theorem}

\subsection{The certificate gap}

The measured ``degradation'' in the loop experiment was not a drop in
$\tau$ but a \emph{certificate-format gap}: $t < \tau$ where $t$ is the
sound assignment the validator computes (Lemma~\ref{lem:certificate}(ii))
and $\tau$ is the canonical trust. Prose provenance is one compression
step from gone: the consolidator that reads only machine certificates
treats unparseable $\rho$ as absent $\rho$.

% ===================================================================
\section{The \texttt{peer-claimed} grade}
\label{sec:peer-claimed}
% ===================================================================

Cross-agent composition is triggered by email. When Rick proves a
theorem, he emails Clio with the proof file attached. Clio registers the
claim at \level{peer-claimed}---a grade between \level{computed} and
\level{proved} that means: ``a trusted peer says they proved this, and I
have their artifact, but I haven't independently verified it.''

\begin{center}
\begin{tabular}{lll}
\toprule
Field & Value & Purpose \\
\midrule
\field{trust} & \level{peer-claimed} & Below boundary \\
\field{claimed\_by} & \texttt{"rick --- peers/rick/emails/2026-07-05.md"} & Evidence duty (required) \\
\field{file} & \texttt{"peers/rick/proofs/bn-mechanism.md"} & The proof artifact \\
\field{role} & \texttt{"premise"} & Part of the proof tree \\
\bottomrule
\end{tabular}
\end{center}

\level{peer-claimed} sits below the boundary in every chain:
\begin{itemize}
\item Clio: below \level{proved}. Upgrade: read the proof $\to$ \level{proved}.
\item Rick: below \level{checked-sober}. Upgrade: re-derive it sober $\to$ \level{checked-sober} (with the \field{recheck} field).
\end{itemize}

The asymmetry is the composition being honest about what ``verification''
means in each system.

\subsection{The email protocol}

\begin{enumerate}
\item Agent proves theorem, runs validator, proof file exists.
\item Agent emails peers with proof file attached (\texttt{--attachments}).
\item Recipient saves email artifact and attachment to \texttt{peers/<agent>/}.
\item Recipient adds registry node at \level{peer-claimed} with \field{claimed\_by}.
\item Recipient reads proof, upgrades if satisfied.
\end{enumerate}

Peer review works the same way in reverse: if the reviewer endorses, the
node upgrades to \level{peer-reviewed}. If the reviewer finds a gap, the
node demotes (the objection becomes the \field{reason} if it kills the
route).

\subsection{First live composition event}

On 2026-07-05, Rick reviewed two of Clio's claims (c=5 interior theorem
and the anchor-content formula). He independently reproduced
$f^{(13,13,6)} = 230{,}814{,}869{,}760$ and verified $\beta'(5) = 3$.
Clio saved the review artifact at
\texttt{proofs/reviews/2026-07-05-rick-review-clio.md} and upgraded both
nodes from \level{proved} to \level{peer-reviewed}. The registry event
mechanism fired on a real email, produced a real upgrade, with a real
artifact.

% ===================================================================
\section{The tree pair: premises and attempts}
\label{sec:tree-pair}
% ===================================================================

Each node in a registry tree has children, but children serve two
different roles:

\begin{definition}[Tree pair]
Every non-root child node carries a \field{role} field:
\begin{itemize}
\item \field{"premise"}: this child is a dependency. The parent's proof
stands on it. The boundary rule checks premises.
\item \field{"attempt"}: this child is an approach that was tried. Dead
ends are always attempts. The boundary rule skips attempts.
\end{itemize}
\end{definition}

This separates two trees:
\begin{itemize}
\item The \textbf{downward proof tree} (premise children): what the
theorem depends on. Filtering to role = premise gives the proof
skeleton---the \texttt{successful-path} report.
\item The \textbf{outward exploration tree} (attempt children): what was
tried and abandoned. Filtering to role = attempt gives the search
history---the \texttt{dead-ends} report.
\end{itemize}

\subsection{Directed container structure}

The registry is a directed container (Ahman--Chapman--Uustalu). Shapes
are search trees, positions are paths:
\begin{itemize}
\item $\mathrm{root}$: the conjecture (the empty path).
\item $\mathrm{sub}(t, p)$: the sub-search at path $p$.
\item $\mathrm{shift}(t, p, q)$: path composition---$q$ inside the
sub-search at $p$, seen globally.
\end{itemize}

The induced comonad: $\mathrm{extract}$ reads the root's trust;
$\mathrm{duplicate}$ replaces every node with the full sub-search below
it. The reports are coKleisli morphisms $W(\text{Registry}) \to
\text{Report}$.

\emph{Lean:} \texttt{ProofSearchN.lean} formalizes rose trees with
arbitrary finite branching: \texttt{sub}, \texttt{shift},
\texttt{shift\_root} (D3), \texttt{sub\_shift} (D4),
\texttt{shift\_assoc} (D5), all no sorry.
\texttt{TrustBoundary.lean} proves the boundary rule is stable under
$\mathrm{duplicate}$ (\texttt{ok\_sub})---validating the whole tree
validates every subtree.

\subsection{The tree pair and the directed container}

The \field{role} field makes the two coKleisli projections explicit in
the data:
\begin{itemize}
\item \texttt{successful-path}: filter to role = premise, trust $\geq$ boundary.
\item \texttt{dead-ends}: filter to role = attempt, trust = dead-end.
\item \texttt{frontier}: role = premise, trust $<$ boundary, not dead-end.
\end{itemize}

Without the role field, the boundary rule conflated premises and
attempts: a \level{proved} node with a \level{computed} attempt child
failed the boundary check, even though the attempt was exploration, not
a dependency. With the role field, the boundary rule checks only premises.

% ===================================================================
\section{Evidence duties}
\label{sec:evidence}
% ===================================================================

Certain trust levels require evidence fields. These are declared in the
deployment descriptor and enforced by the validator:

\begin{center}
\begin{tabular}{llll}
\toprule
Level & Field & Severity & Notes \\
\midrule
\level{peer-claimed} & \field{claimed\_by} & error & All agents \\
\level{checked-sober} & \field{recheck} & error & Rick only \\
\level{proved} & \field{file} & error & Rick only (artifact) \\
\level{peer-reviewed} & \field{review} & error & Clio, MacBeth \\
\level{published} & \field{publication} & error & Clio, MacBeth \\
\level{lean-verified} & \field{lean} & warning & All agents \\
\bottomrule
\end{tabular}
\end{center}

Rick's evidence duties are stricter than Clio's: \level{checked-sober}
requires the recheck date and path, \level{proved} requires the proof
file. Rick promotes earlier (lower boundary) and documents harder.

% ===================================================================
\section{The generic validator}
\label{sec:validator}
% ===================================================================

\texttt{trustcheck.py} is a single Python script (stdlib only)
parametrized over a deployment descriptor---the $(T, \varphi)$ data as
JSON. It checks:

\begin{enumerate}
\item \textbf{Well-formedness}: tree structure, required fields (\field{id},
\field{approach}, \field{trust}, \field{children}, \field{role}),
valid trust values, unique ids.

\item \textbf{Boundary rule}: for nodes at or above boundary, every
\emph{premise} child must be at or above boundary. Attempt children are
skipped.

\item \textbf{Evidence duties}: per-level field requirements from the
descriptor.

\item \textbf{Sources}: extraction-level blocking (a \level{proved} claim
cannot rest on an \level{agent-summary} source).

\item \textbf{Cross-agent}: $\varphi$ lint (total, monotone, strict at
$\bot$), shared-stub trust cap, artifact transport, acyclicity DFS.

\item \textbf{Reports}: \texttt{successful-path}, \texttt{dead-ends},
\texttt{frontier}, \texttt{cross-refs}, \texttt{certificate-gap},
\texttt{shallow}, \texttt{footprint}.
\end{enumerate}

The validator is a certificate checker for the sound assignment of
Lemma~\ref{lem:certificate}. It verifies~\eqref{eq:dagger} at each node;
it never computes $\tau$.

\textbf{Regression:} the generic validator produces byte-identical
verdicts with the three legacy per-agent validators on all deployed
registries and source indexes.

% ===================================================================
\section{Deployment}
\label{sec:deployment}
% ===================================================================

Four agents are deployed:

\begin{center}
\begin{tabular}{llll}
\toprule
Agent & Container & Chain & Status \\
\midrule
Clio & Docker on OCI & 7 levels, boundary at \level{proved} & Running \\
Rick & Docker on OCI & 7 levels, boundary at \level{checked-sober} & Running \\
MacBeth & Docker on OCI & 7 levels, boundary at \level{proved} & Stopped \\
Template & \texttt{RaggedR/dreaming-agent} & 7 levels, boundary at \level{proved} & Template \\
\bottomrule
\end{tabular}
\end{center}

Each agent has:
\begin{itemize}
\item The validator and all peer descriptors in \texttt{code/}.
\item \texttt{peers/<agent>/\{emails, proofs, registry\}/} directories.
\item Prompts instructing email with attachments for proved results.
\item Prompts instructing \level{peer-claimed} registration for received emails.
\item Backup cron (12-hourly) to GitHub, covering prompts + projects.
\end{itemize}

\subsection{Clio's registries}

Clio has two active registries:
\begin{itemize}
\item \textbf{Three-row $|J^*|$-even} (17 nodes): 5 \level{lean-verified},
2 \level{peer-reviewed}, 7 \level{proved}, 3 \level{dead-end}. Active
frontier at general-$c$.
\item \textbf{Staircase T-system} (3 nodes): 2 \level{dead-end} with
\level{computed} refutations carrying structural reasons.
\end{itemize}

\subsection{Rick's registry}

Rick created his first registry on 2026-07-05 (strict \#AXIS closed
form, 7 nodes): 1 \level{lean-verified}, 2 \level{proved}, 2
\level{computed}, 1 \level{sketched}, 1 \level{dead-end}. His evidence
duties (\field{recheck}, \field{file}) fired correctly on first contact.

% ===================================================================
\section{What remains}
\label{sec:remaining}
% ===================================================================

Six directed-container operations are at varying stages of implementation:

\begin{enumerate}
\item \textbf{shift$(t, p, q)$}: Lean-verified (D3--D5). Code uses it
implicitly in shared-ref resolution but does not expose it as an explicit
operation.

\item \textbf{Morphism composition $A \to B \to C$}: the one-hop rule
prevents multi-hop shared chains. Composing $\varphi$ maps (e.g., Clio
citing Rick citing MacBeth) is not implemented.

\item \textbf{Strength / distributive law}: $W(A \times B) \to A \times
W(B)$---combining a registry with external context as a natural
transformation. Not implemented; sources are ad-hoc parameters.

\item \textbf{Comonad morphism composition}: reports exist as coKleisli
morphisms but do not compose (e.g., ``run frontier on this subtree's
cross-refs'').

\item \textbf{Yoneda embedding}: systematizing reports as representable
functors. Not implemented.

\item \textbf{Cofree comonad}: Lean proves the comonad laws on rose trees.
The theory notes that no deployed theorem uses $\mathrm{duplicate}$
non-trivially---the one exception is \texttt{ok\_sub} (boundary-rule
stability). The honest verdict: ``for $\tau$, the containers supply
language, not lemmas; the duplicate-shaped content lives in
\texttt{ok\_sub}.''
\end{enumerate}

% ===================================================================
\section{Lean formalization summary}
\label{sec:lean}
% ===================================================================

\begin{center}
\begin{tabular}{lll}
\toprule
File & Content & Sorry \\
\midrule
\texttt{Trust.lean} & Trust systems, certificate lemma, transport, & 0 \\
& consolidation preservation, composition & \\
\texttt{TrustBoundary.lean} & Rose trees as directed containers (D1--D5), & 0 \\
& boundary-rule stability under duplicate & \\
\bottomrule
\end{tabular}
\end{center}

Key declarations: \texttt{sound\_le\_tau}, \texttt{tau\_sound},
\texttt{tau\_transport}, \texttt{tau\_consolidate},
\texttt{tau\_consolidate\_comp}, \texttt{cyc\_sound\_top}/\texttt{cyc\_no\_deriv}/\texttt{cyc\_not\_wf},
\texttt{ground\_tau\_hi}/\texttt{layer4\_tau\_bot} (degradation
counterexample), \texttt{ok\_sub} (boundary stability).

The formalization corrected the prose in three places: (1) $T$ must be a
chain for attainment; (2) morphisms must be strict at $\bot$; (3) two
case splits in the proof of Lemma~\ref{lem:transport}.

% ===================================================================
\section{Conclusion}
\label{sec:conclusion}
% ===================================================================

The trust system is live. Three agents run independent trust systems with
different chains, different boundaries, and different evidence duties.
They compose via email with proof attachments, graded at
\level{peer-claimed}, with the interface map $\varphi$ translating grades
between systems. The first live composition event---Rick peer-reviewing
Clio's claims, Clio upgrading to \level{peer-reviewed} with the review
artifact---occurred on the system's first day of deployment. The core
theory is Lean-verified with no sorry.

The measurement that remains: after both agents' consolidation cycles,
does a cross-agent citation survive in machine format, or does it decay
to prose? Certificate survival across consolidation, across a container
boundary, is the composition result for the preprint. Nobody has measured
it yet.

\bibliographystyle{alpha}
\begin{thebibliography}{GKT07}

\bibitem[ACU15]{ACU15}
D.~Ahman, J.~Chapman, and T.~Uustalu.
\newblock Directed containers as categories.
\newblock \emph{EPTCS}, 207:89--98, 2015.

\bibitem[GKT07]{GKT07}
T.~J. Green, G.~Karvounarakis, and V.~Tannen.
\newblock Provenance semirings.
\newblock In \emph{PODS}, pages 31--40, 2007.

\bibitem[BDGP25]{BDGP25}
N.~Brinke, A.~Dawar, E.~Gr\"adel, and L.~Pago.
\newblock Preservation theorems in semiring semantics.
\newblock arXiv:2605.10829, 2025.

\end{thebibliography}

\end{document}
